Figure 1. Geometry used in cross-mesh integration. (a) One target triangle crosses several source triangles. The darker polygon is its intersection with one source triangle and gives one region for integrating the Galerkin right-hand side. (b) The six triangles and quadrilaterals around vertex $k$ each contribute a shaded piece to its median-dual control volume $C_k^M$. The blue boundary joins edge midpoints to cell centers. Dotted lines in one quadrilateral show the triangles used for reconstruction. Only the original mesh vertices carry solver unknowns. Both panels are illustrative mesh patches.

Figure 2. Exact density before transfer in the two static tests, shown on the same scale. The smooth field varies continuously. The circular inclusion has density 3.6~kg/m$^3$ inside and 1.2~kg/m$^3$ outside.

Figure 3. Smooth-field transfer accuracy on structured quadrilateral (left) and unstructured mixed (right) target meshes. The vertical axis shows the relative nodal density error $e_{\rho,2}^{\mathrm{node}}$; $h_{\mathrm{eff}}$ is the effective spacing derived from the vertex count and decreases to the right. Gray lines show first-, second-, and third-order reference slopes. Curves of equal order nearly coincide; Table~\ref{tab:observed-orders} gives the observed orders.

Figure 4. Density at $t=0.399$ after 398 transfers, for $s=2$ and mesh sequence 1. Each numerical profile averages eight lines of constant $y$. Shading marks $[2.40,2.46]$, where the smeared contact affects the NN1 and BAR2 profiles. The shock width is measured in $[2.55,2.85]$.

Figure 5. Shock broadening and mass change at $t=0.399$ after 398 transfers with $s=2$. The three symbols for each method are individual mesh sequences, offset horizontally. (a) Shock width relative to the fixed-mesh run, in reference cells. (b) Magnitudes of the signed sum of mass changes during transfer and of the final mass drift, both divided by the initial mass.

Figure 6. First-step lift response to the same error field imposed in one stored state, divided by the response when it is imposed in both. Dashed lines are the ratios $4/3$ and $-1/3$ given by the BDF2 weights in Eq.~\eqref{eq:naca-history-ratios}. The differences are taken from the reference run on the target Euler airfoil mesh.

Figure 7. Pressure response to a single Euler remeshing event in the annulus $0.20c\le r\le0.25c$, where $r$ is distance from the airfoil wall and $c$ is chord. (a) Shell RMS pressure differences from the fixed-mesh reference, normalized separately by each shell's maximum over the 400 post-event samples; shell half-width is $0.004c$. Disturbances reach shells closer to the wall later. (b) First crossings of 50\% of each shell's maximum for GP2 and CV2, restricted to travel distances of at least $0.08c$. The dotted line gives the travel time at the freestream sound speed. Each method is represented by one event.

Figure 8. Time-step sensitivity of drag with SA over three remeshing events for mesh sequence 1. Panel (a) shows the GP2 drag difference $\Delta C_D$ from the fixed-mesh reference at the same time step. Panel (b) shows the NN2--GP2 drag difference, $\delta C_D=C_{D,\mathrm{NN2}}-C_{D,\mathrm{GP2}}$. All curves use common physical sample times; the horizontal axis gives time divided by the nominal step $\Delta t_0$. Mesh perturbations are applied at matching physical times; dotted vertical lines mark the first sample after each of the three remeshing events.
